\section{Domain of Hopes and Dreams}

\subsection*{Deskripsi Soal}

Dalam suatu catatan simposium, Marisa menemukan rekaman mengenai keberadaan
sebuah dimensi buatan yang terisolasi dari dunia luar. Batas-batas menuju
wilayah rahasia ini dikunci menggunakan segel jimat yang dienkripsi dengan
\textbf{Playfair \textit{Cipher}} agar tidak dapat diakses secara sembarangan.

Tugas kami adalah memulihkan naskah petunjuk akses tersebut. Susun kembali
matriks kunci dengan menganalisis pola kemunculan pasangan huruf (\textit{bigram}
maupun \textit{n-gram} lainnya), distribusi karakter, serta tebakan kata-kata
umum dalam \textbf{Bahasa Inggris}. Setelah matriks kunci berhasil
direkonstruksi, dekripsikan seluruh dokumen rahasia tersebut.

\medskip
\noindent\textbf{Berkas lampiran:}
\begin{itemize}[nosep]
  \item \textit{Attachment}: \url{https://drive.google.com/file/d/1ODljf85XJXS41PSSqrfaqUOQywTuneuX/view?usp=sharing}
  \item \texttt{md5sum}: \texttt{b6f21903cbdb62ca899d6bb2694814b1}
  \item \texttt{sha256sum}: \texttt{e4bb0e2a2ac8425ef0133b78dd0103308f781d2afdfcc46712b2a23efb7ee597}
\end{itemize}

\subsection*{Berkas \textit{Ciphertext}}

\ciphertextfile{../../src/23/f.txt}

\subsection*{Langkah Penyelesaian}

\textbf{Playfair \textit{Cipher}} sangat berkaitan dengan penggunaan bigram (karena aturan enkripsinya berpusat
pada pemetaan bigram di sebuah matriks kunci $5 \times 5$). Kelemahan \textit{cipher} ini:

\begin{itemize}
  \item Rentan terhadap analisis frekuensi \textbf{bigram}, jika teks \textit{cipher} cukup panjang.
  \item Bigram dan kebalikannya akan menghasilkan \textit{cipher} yang saling berkaitan: jika bigram
  plaintext $XY$ (baik hubungan searah baris, kolom, maupun \textit{rectangle}) dienkripsi menjadi
  \textit{cipher} $PQ$, maka bigram kebalikannya $YX$ akan dienkripsi menjadi $QP$, kebalikan dari
  \textit{cipher} sebelumnya.
\end{itemize}

Penyelesaian soal ini dibagi menjadi dua tahap. Tahap pertama adalah \textbf{analisis frekuensi bigram
secara manual} untuk menambatkan sebagian huruf pada matriks kunci. Tahap kedua adalah \textbf{pencarian
otomatis dengan \textit{Simulated Annealing} (SA)} untuk melengkapi sel-sel yang tersisa, karena
melanjutkan penebakan manual sel demi sel terlalu lambat.

\subsubsection*{Tahap 1: Analisis Frekuensi Bigram}

\begin{enumerate}
  \item Hitung frekuensi bigram \textbf{non-\textit{overlapping}} pada \textit{ciphertext}
  (\texttt{Solver.freq\_bigram} pada \texttt{src/23/solver.py}). Perhitungan ini dilakukan 
  pada \textit{ciphertext} \textbf{sebelum} dilakukan substitusi
  apa pun, sehingga hasilnya tidak berubah walau hipotesis matriks kunci diubah-ubah. Terdapat
  \textbf{392} bigram unik, dengan sepuluh besar sebagai berikut.

  \begin{center}
    \small
    \begin{tabular}{|c|c|c||c|c|}
      \hline
      \multicolumn{5}{|c|}{Bigram (non-\textit{overlapping})} \\
      \hline
      Peringkat & Cipher & Freq (\%) & Referensi & Freq (\%) \\
      \hline
      1  & BF & 3.24 & TH & 3.56 \\
      2  & UA & 1.95 & HE & 3.07 \\
      3  & GB & 1.87 & IN & 2.43 \\
      4  & VT & 1.62 & ER & 2.05 \\
      5  & TV & 1.60 & AN & 1.99 \\
      6  & YA & 1.45 & RE & 1.85 \\
      7  & IA & 1.34 & ON & 1.76 \\
      8  & DS & 1.22 & AT & 1.49 \\
      9  & ID & 1.22 & EN & 1.45 \\
      10 & QL & 1.21 & ND & 1.35 \\
      \hline
    \end{tabular}
  \end{center}

  \item Susun hipotesis awal dari empat pencocokan berikut:

  \begin{center}
    \small
    \begin{tabular}{|c|c|c|c|}
      \hline
      Cipher & Freq (\%) & Dugaan \textit{Plain} & Dasar Dugaan \\
      \hline
      \texttt{BF} & 3.24 & \textbf{TH} & bigram tersering pada kedua sisi \\
      \texttt{GB} & 1.87 & \textbf{HE} & peringkat 3 cipher vs bigram Inggris tersering kedua \\
      \texttt{VT} & 1.62 & \textbf{RE} & pasangan resiprokal dengan \texttt{TV} \\
      \texttt{TV} & 1.60 & \textbf{ER} & pasangan resiprokal dengan \texttt{VT} \\
      \hline
    \end{tabular}
  \end{center}

  \texttt{UA} belum dicocokkan karena cukup sulit memetakan huruf \textit{plain} yang beririsan. 
  Melanjutkan ke \texttt{VT}/\texttt{TV}, keduanya muncul dengan frekuensi yang mirip, sehingga hipotesis
  selanjutnya adalah memetakan \texttt{VT}/\texttt{TV} ke \textbf{RE}/\textbf{ER}, karena memang 
  keduanya merupakan pasangan bolak-balik yang sama-sama sering muncul dalam Bahasa Inggris 
  (2.05\% dan 1.85\%). Pencocokan \texttt{GB} $\to$ \textbf{HE} juga
  terverifikasi lewat kemunculan \texttt{BG} $\to$ \textbf{EH}.

  \item Turunkan hubungan posisi antarhuruf dari keempat hipotesis tersebut. Ketiganya ternyata
  merupakan hubungan \textit{rectangle}, sehingga masing-masing mengunci baris dan kolom relatif
  antarhuruf:

  \begin{itemize}[nosep]
    \item \texttt{BF} $\to$ \textbf{TH}: \texttt{B} sebaris dengan \texttt{T}, dan \texttt{F} sebaris dengan \texttt{H}.
    \item \texttt{GB} $\to$ \textbf{HE}: \texttt{G} sebaris dengan \texttt{H}, dan \texttt{B} sebaris dengan \texttt{E}.
    \item \texttt{TV} $\to$ \textbf{ER}: \texttt{T} sebaris dengan \texttt{E}, dan \texttt{V} sebaris dengan \texttt{R}.
  \end{itemize}

  Gabungan ketiganya membentuk blok \texttt{T}-\texttt{E}-\texttt{B} pada satu baris,
  \texttt{F}-\texttt{G}-\texttt{H} pada baris lain, serta \texttt{R}-\texttt{V} pada baris ketiga, dengan
  kolom yang saling bersesuaian. Penambahan \texttt{TF} $\to$ \textbf{IT} kemudian menempatkan
  \texttt{I}. Total \textbf{9 dari 25} sel berhasil ditambatkan secara manual, dan inilah isi
  \texttt{GRID} pada \texttt{src/23/solver.py}.

  Catatan: hipotesis sempat keliru ketika \texttt{R} dan \texttt{V} diletakkan pada baris
  yang sama dengan \texttt{T}, \texttt{E}, dan \texttt{B}. Susunan itu memaksa \texttt{BE}
  (0.88\%) terbaca sebagai \textbf{VB}, sebuah bigram yang praktis tidak muncul dalam Bahasa
  Inggris. Kontradiksi inilah yang menjadi dasar penolakan hipotesis tersebut.
\end{enumerate}

Hasil tahap ini adalah hipotesis matriks kunci berikut, dengan sel kosong menandakan posisi yang belum
dapat disimpulkan dari analisis bigram:

\begin{center}
  \begin{tabular}{|c|c|c|c|c|}
    \hline
      &   &   &   &   \\
    \hline
    I &   &   &   &   \\
    \hline
    T & E & B &   &   \\
    \hline
    F & G & H &   &   \\
    \hline
    R & V &   &   &   \\
    \hline
  \end{tabular}
\end{center}

\subsubsection*{Tahap 2: \textit{Simulated Annealing}}

Melanjutkan tahap 1 secara manual untuk 16 sel sisanya tidak efisien. Seringkali kesalahan posisi satu huruf
baru ditemukan beberapa langkah kemudian. Karena itu, sisa pencarian diserahkan pada \textit{Simulated Annealing}.

\begin{enumerate}
  \item Hasil dekripsi dinilai dengan \textbf{fungsi objektif} berupa \textbf{log-probabilitas quadgram}
  (\texttt{Solver.score}), yaitu
  \[
    S = \sum_{i} \log_{10} P(q_i),
  \]
  dengan $q_i$ adalah quadgram ke-$i$ (jendela 4 huruf yang saling tumpang-tindih) pada hasil dekripsi.
  Nilai $P$ diambil dari tabel \texttt{QUADGRAM} pada \texttt{src/frequencies.py} (\textit{quadgram} yang tidak
  terdaftar diberi nilai dasar $0.001$ sebagai penalti). \textit{Quadgram}dipilih karena menilai kelayakan
  rangkaian huruf yang lebih panjang daripada bigram, sehingga lebih tegas membedakan teks Inggris dari
  teks acak.

  \item Satu langkah dilakukan dengan menukar isi dua sel bebas secara acak
  (\texttt{Solver.swap}). Penukaran antarbaris atau antarkolom penuh hanya diaktifkan bila seluruh sel
  bebas, karena operasi tersebut akan menggeser sel-sel yang sudah dikunci.

  \item Solusi tetangga dengan skor lebih tinggi selalu diterima. Kemudian, solusi yang lebih buruk tetap 
  diterima dengan peluang $\exp(\Delta S / T)$, sehingga pencarian dapat keluar
  dari optimum lokal. Suhu awal $T = 20$ diturunkan bertahap sebesar $0.5$ hingga mencapai $0$, dengan
  $2500$ percobaan pada tiap tingkat suhu (total $\pm 100.000$ evaluasi per \textit{restart}). Proses ini
  diulang sebanyak 2 \textit{restart}, dan skor terbaiklah yang akan dipakai.

  \item Selama pencarian, penilaian hanya dilakukan pada \textbf{2400 karakter pertama} \textit{ciphertext}. 
  Ukuran ini sudah cukup untuk membedakan kunci benar dari kunci salah,
  namun memangkas waktu secara signifikan. Dekripsi penuh baru dijalankan sekali di akhir
  memakai kunci terbaik.
\end{enumerate}

\subsubsection*{Hasil}

Pencarian konvergen ke skor $-6877.2$ dengan waktu eksekusi sekitar \textbf{2 menit 30 detik}, dan
menghasilkan pemetaan lengkap untuk seluruh \textbf{392} bigram unik. Pengujian berulang memberi kunci
yang sama secara konsisten.

Perlu dicatat bahwa penguncian sembilan sel hasil tahap 1 sangat menentukan. Ketika seluruh 25 sel
dibiarkan bebas, pencarian dengan jumlah \textit{restart} yang sama kerap berhenti pada optimum lokal
(misalnya pada skor $-7098.3$) dengan hasil dekripsi yang tidak terbaca. Dengan kata lain, analisis manual
dan pencarian otomatis saling melengkapi: analisis manual memperkecil ruang pencarian dari 25 sel menjadi
16 sel, sementara SA menyelesaikan bagian sisanya.

\subsection*{Artefak}

Skrip \textit{solver} yang digunakan untuk menghitung frekuensi bigram \textit{ciphertext}, menjalankan
\textit{Simulated Annealing} untuk melengkapi matriks kunci, menurunkan pemetaan bigram
\textit{cipher}$\to$\textit{plain} berdasarkan aturan Playfair, mendekripsi \textit{ciphertext}, serta
mencatat setiap hasil dekripsi ke \texttt{log/plain/} dan metadatanya (pemetaan, matriks kunci, frekuensi
bigram, dan skor akhir) ke \texttt{log/metadata/}.

\lstinputlisting[language=Python, basicstyle=\ttfamily\small, frame=single,
  title=\texttt{src/23/solver.py}]{../../src/23/solver.py}

\subsection*{\textit{Plaintext} Hasil Dekripsi}

\ciphertextfile{../../src/23/log/plain/2026-09-15_171115.txt}

\subsection*{\textit{Plaintext} dengan Format Asli}

\begin{tcolorbox}[breakable, colback=white, colframe=black!80,
  boxrule=0.6pt, arc=0mm, left=3mm, right=3mm, top=2mm, bottom=2mm]
\ttfamily\small\raggedright
Part I\par
What Change Will the New Powers Bring to Gensokyo?\par
\medskip
I\textquotesingle{}d like to thank everyone for taking the time out of their busy lives to participate. I wish to collect your valued views and report them for the benefit of all Gensokyo. With that in mind, I\textquotesingle{}d like to turn it over to our symposium\textquotesingle{}s moderator, human magician and one of leading experts on youkai.\par
\medskip
Ah, um, I am the moderator. \textquotesingle{}Cause... Uh, I mean, because I believed the symposium to be presented here may be interesting, I have accepted the honor of this role. And, umm...\par
\medskip
Just speak as you normally do. Hearing you talk like this is disgusting.\par
\medskip
Okay, then. Let\textquotesingle{}s start with the introductions. Take it away.\par
\medskip
I manage the shrine on Youkai Mountain. I usually think of nothing but Shinto, but I am interested to hear from those of other religions today.\par
\medskip
I am the chief priest of the temple on the outskirts of the village. Although I\textquotesingle{}m the chief priest, I\textquotesingle{}m still a novice yet to attain enlightenment, but please consider my views today.\par
\medskip
Um, presently, I\textquotesingle{}ve renounced my humanity and am a hermit. I only awoke recently, so I\textquotesingle{}m still half-{}asleep, but I felt that I might deepen my understanding of Gensokyo by attending today, so I was determined to do so. Speaking of which, the residents of Gensokyo seem to be on good terms, if they can have a symposium like this...\par
\medskip
Whoa, this is just the first one. But we\textquotesingle{}re always up fer takin\textquotesingle{} on new challenges, though.\par
\medskip
Is that so...\par
\medskip
[What Change Will the New Powers Bring to Gensokyo?]\par
\medskip
So we got here a Shinto god, a Buddhist priest, and a Taoist hermit. Since we\textquotesingle{}ve pretty much gathered up a bunch of business rivals, this oughta be good. A three-{}way talk between three religious figures. So, what are y\textquotesingle{}all going to talk about?\par
\medskip
Aren\textquotesingle{}t you, our moderator, supposed to decide that?\par
\medskip
Oh, yeah. Okay, then, the first theme is \textquotedbl{}What change will the new powers bring to Gensokyo?\textquotedbl{} Just throw whatever ya want out there.\par
\medskip
As you already are. (laugh)\par
\medskip
The \textquotedbl{}new powers\textquotedbl{} refers to all of us, right?\par
\medskip
Ah, right. Just to clarify a little, y\textquotesingle{}all three are the newcomers. The hidden theme of the symposium is that ya might be hostile elements here to disrupt our peaceful Gensokyo or somethin\textquotesingle{}.\par
\medskip
And here I thought you were starting to get used to me...\par
\medskip
Don\textquotesingle{}t matter how many years go by, newcomers will always be newcomers. From the old timers\textquotesingle{} perspective, anyway.\par
\medskip
Well then, as the real newcomer, I shall start.\par
\medskip
Go for it.\par
\medskip
My first impression of Gensokyo, besides the fact that there are youkai who cause harm to mankind, is that the minds of humans are capricious and easily swayed. The humans of the village spend their days in mental vagrancy, failing to improve their lot in life, and the youkai are allowed to take advantage of their capriciousness to their heart\textquotesingle{}s content. Frankly, it\textquotesingle{}s a miracle that society hasn\textquotesingle{}t collapsed yet, and I believe the current balance is dangerously unstable.\par
\medskip
Uh, I guess...\par
\medskip
Why does Gensokyo have no administrator? Society cannot function with all humans possessing the same social standing and rank.\par
\medskip
An administrator?\par
\medskip
That\textquotesingle{}s right. Humans need leaders to guide them towards a hope-{}filled future. If no one else is qualified then I would certainly...\par
\medskip
I wonder what kind of future lies in store for the humans under the rule of someone who is no longer human.\par
\medskip
\textquotedbl{}Rule\textquotedbl{} is such an ugly word. It would just be a stopgap measure until someone from among the humans is able to imitate and overturn me.\par
\medskip
May I say a word? You say that there is no administrator in Gensokyo, but I don\textquotesingle{}t think that\textquotesingle{}s quite true.\par
\medskip
What do ya mean?\par
\medskip
Gensokyo\textquotesingle{}s current form cannot be maintained without the youkai\textquotesingle{}s power. Gensokyo was created by youkai to save themselves, after all. The humans in Gensokyo exist for no further purpose than to preserve the youkai. The humans are being ruled by the youkai.\par
\medskip
I agree. The youkai keep watch on each other to prevent other youkai from running rampant and ruling over the humans, and to make sure the humans don\textquotesingle{}t awaken to their actual condition and turn down the wrong path. Therefore, the fact that a ruler hasn\textquotesingle{}t emerged from among the humans is in Gensokyo\textquotesingle{}s best interest.\par
\medskip
So that\textquotesingle{}s how it is. But if that\textquotesingle{}s the case, then the poor humans here are pitiable creatures...\par
\medskip
Think of it this way: are the animals in a zoo unfortunate?\par
\medskip
And a zoo is?\par
\medskip
An institution where endangered animals are protected and humans can observe them and learn about them. Taking a naturally free bird and locking it up in small cage... when you look at it from that perspective, it seems quite unfortunate. But now, it doesn\textquotesingle{}t have to worry about predators, or running out of food, and there will always be plenty of humans for it to look at, so it will never get bored. If you think of it like that, aren\textquotesingle{}t they fortunate?\par
\medskip
So Gensokyo\textquotesingle{}s humans are animals in a zoo? What happened to human dignity?\par
\medskip
Dignity, huh? I don\textquotesingle{}t think that\textquotesingle{}s something you receive from others, though... What do you think?\par
\medskip
Hmm? Me? What\textquotesingle{}s dignity?\par
\medskip
Like pride, or self-{}respect.\par
\medskip
Of course, I\textquotesingle{}m proud of everything I do, but even if I were forced inta somethin\textquotesingle{}, I\textquotesingle{}d be perfectly fine with that so long as I were havin\textquotesingle{} fun.\par
\medskip
I see... Then please forget about what I said.\par
\medskip
Wait, what were ya talkin\textquotesingle{} about?\par
\medskip
(Laugh)\par
\medskip
[The Reason for Coming to Gensokyo]\par
\medskip
Anyway, I was havin\textquotesingle{} a hard time followin\textquotesingle{} that. Honestly, I don\textquotesingle{}t give a flip about the humans of the village. Talk about yerselves. The theme was \textquotedbl{}what are you bringing to Gensokyo\textquotedbl{}, after all. So first of all, why\textquotesingle{}d ya move to Gensokyo anyway?\par
\medskip
I came to Gensokyo to seek new faith. In the outside world, shrines have become mere tourist traps or called nonsensical things like \textquotedbl{}power spots\textquotedbl{}, so places for humans to offer their faith to the gods have disappeared. Gods are beings that will cease to exist without humans believing in them.\par
\medskip
So ya say, but I don\textquotesingle{}t think ya\textquotesingle{}ve gotten much faith in Gensokyo, either.\par
\medskip
Just believing that we\textquotesingle{}re real is enough. In the outside world, they\textquotesingle{}ve stopped believing that gods even exist. Indeed, everything really changed once we came to Gensokyo. Here, humans believe that gods exist as an everyday fact.\par
\medskip
W-{}well, even if I believed in ya or not, yer right here, ain\textquotesingle{}t ya?\par
\medskip
That\textquotesingle{}s not good enough for the outside world. Even if you perform a miracle right before their eyes, they\textquotesingle{}ll convince themselves it\textquotesingle{}s an illusion and absolutely refuse to accept that it actually happened. Humans are easily swayed creatures, so that\textquotesingle{}s the latest fad.\par
\medskip
A fad?\par
\medskip
Yes, an entirely intentional fad. But, you know, the best part is that even while they don\textquotesingle{}t think that gods or youkai exist, they still believe in aliens.\par
\medskip
Oh? Never heard about that.\par
\medskip
The moment that humans realized that the vastness of space was beyond their ability to control, they started to convince themselves that there must be aliens somewhere out there, even though nobody\textquotesingle{}s confirmed it yet. Rationalism suddenly turns to mysticism when it goes too far. Funny, isn\textquotesingle{}t it?\par
\medskip
I have been curious about this for a while now, but your shrine has multiple divine presences, am I right? Why is that?\par
\medskip
Technically, we\textquotesingle{}ve got multiple, but never mind that for now. Our shrine\textquotesingle{}s been through quite a bit... Well, just let me say the companion deity is a good partner.\par
\medskip
Ooh, I know this one. The companion deity\textquotesingle{}s the one who really rules the shrine, right? And yer just the one runnin\textquotesingle{} stuff out in front.\par
\medskip
Ahem. Although that\textquotesingle{}s not quite how it works, it\textquotesingle{}s not wrong that the shrine belongs to the companion deity. But the companion deity doesn\textquotesingle{}t have much interest in gathering faith, you see. So I have no choice but to gather faith for us by myself. No choice.\par
\medskip
Speaking of which, you are from that region of Shinano, are you not?\par
\medskip
I\textquotesingle{}m not, but the others are. I lived there for a long time, too, though.\par
\medskip
Actually, I\textquotesingle{}m also from Shinano. That brings back so many memories.\par
\medskip
Oh, really?\par
\medskip
Whenever I came there, the lake was always full of such clean water. It was such a beautiful place.\par
\medskip
When was this?\par
\medskip
I\textquotesingle{}m not sure. Maybe about a thousand years ago.\par
\medskip
Nowadays, it\textquotesingle{}s made a complete about-{}face and become a rather cold city... Although I\textquotesingle{}m pretty proud of its excellent factories and the precision machinery they make. There are also some nice hot springs by the train station....\par
\medskip
[The Construction of the Temple]\par
\medskip
Whoa, let\textquotesingle{}s not get caught up in talk of some countryside town no one here knows about. Next up is you. Why\textquotesingle{}d you come to Gensokyo?\par
\medskip
I am just a simple monk in training, but because of a minor disagreement, I was betrayed by humans and sealed away... However, some of my old friends who missed me came to my rescue. And that place just happened to be Gensokyo.\par
\medskip
Now, wait. Does that mean you\textquotesingle{}re here by coincidence?\par
\medskip
Not coincidence, but rather inevitability. Were they not in Gensokyo, my friends would have been unable to maintain their strength.\par
\medskip
You built your temple directly above where I was sleeping. Was that also inevitable?\par
\medskip
What might you be talking about?\par
\medskip
Oh, yeah. The grave was under the temple, wasn\textquotesingle{}t it?\par
\medskip
You built a temple for Buddhism, a religion I used for political measures, on top of it, and I heard it was not just any temple, but a youkai temple full of youkai monks. A temple filled with such detestable, evil power...\par
\medskip
How rude. While it may be filled with youkai, those youkai are more sincere than humans. Just as humans can become hermits and celestials, youkai can become gods and Buddhas. My own power is also as vast as it is thanks to youkai.\par
\medskip
So you\textquotesingle{}re a monk possessed by evil then. A monk who has fallen unthinkably far... or so I\textquotesingle{}d like to say, but as someone who used Buddhism for political ends, I suppose I do not have the right to do so.\par
\medskip
Well, to be honest, we built the temple there because our treasure-{}hunting companion reported that \textquotedbl{}there\textquotesingle{}s something incredible buried down there!\textquotedbl{} and so we surveyed it. When it turned out to be something that seemed more dangerous than what we were expecting, we thought we should do our best to hide it.\par
\medskip
So I was that dangerous-{}seeming thing, then?\par
\medskip
While I don\textquotesingle{}t think my intuition was wrong, I regrettably lacked the power.\par
\medskip
Meaning you couldn\textquotesingle{}t keep me sealed, unfortunately for our nameless so-{}called monk who thought to keep me shut off forever.\par
\medskip
[The Revival of the Saint]\par
\medskip
You can fight later. So, last up is the final story.\par
\medskip
I put myself to sleep with a spell. I had planned to wake up as a hermit and enjoy the world as I desired, but I awoke once I was aware that I was no longer in the ordinary world. It was then that I came to learn I was here in Gensokyo. I\textquotesingle{}m in the process of carefully studying this world, and today will be particularly informative. I still haven\textquotesingle{}t figured out the reason I\textquotesingle{}m here, though...\par
\medskip
The reason you\textquotesingle{}re here is simple. In the same way that youkai or gods like me end up here after their existence is renounced in the outside world, you\textquotesingle{}ve also ceased to exist.\par
\medskip
I\textquotesingle{}ve ceased to exist? But, aren\textquotesingle{}t I here right now...?\par
\medskip
Common sense and uncommon sense; in other words, reality and fantasy. The outside world and Gensokyo are separated along that same line. Hence, things that have disappeared, been forgotten, or have had their existence renounced in the outside world can be found in Gensokyo.\par
\medskip
So that means I was...\par
\medskip
Completely forgotten.\par
\medskip
So I\textquotesingle{}ve changed from a human to a god-{}like existence, then?\par
\medskip
That\textquotesingle{}s correct. Well, as far as I know you were already regarded as a god in Shinto.\par
\medskip
Oh, yes. I forgot to mention this but I don\textquotesingle{}t normally live in Gensokyo.\par
\medskip
Huh?\par
\medskip
I live in a space I created myself. Living somewhere where nobody can bother them is the essence of what a hermit is.\par
\medskip
Now thatcha mention it, another hermit I know also lives in a hard-{}to-{}reach place. Howd\textquotesingle{}ya do somethin\textquotesingle{} like that anyway?\par
\medskip
It\textquotesingle{}s simple, if you study enough of the hermetic arts. Would you like to become my disciple?\par
\medskip
Uh, lemme think about it.\par
\medskip
You could become immortal.\par
\medskip
Ooh.\par
\medskip
Oh, proselytizing now, are we? Don\textquotesingle{}t think I\textquotesingle{}ll take that lying down. Switch now and I\textquotesingle{}ll throw in excellent luck on all your fortunes for two months.\par
\medskip
Then allow me as well. You are a youkai-{}like human, after all.\par
\medskip
You certainly are popular. (laugh)\par
\medskip
A-{}Anyway, I guess that wraps up the self-{}introductions. So what I want to ask y\textquotesingle{}all now is whaddya want to do in Gensokyo from now on?\par
\medskip
I think I\textquotesingle{}d like to live a carefree life without being disturbed.\par
\medskip
Izzat so?\par
\medskip
Transcending worldly matters, to live a refined life in a world with no worries or anything at all, that\textquotesingle{}s my ideal.\par
\medskip
You might have been thinking you\textquotesingle{}ve managed to become immortal and ageless by studying hermit techniques, but unfortunately, immortals in Gensokyo have a natural enemy.\par
\medskip
You mean shinigami?\par
\medskip
Yes, they will periodically come to kill you.\par
\medskip
I believe those battles will be a form of entertainment. A nice, leisurely struggle.\par
\medskip
Incidentally, I am not interested in immortality. Well, to be honest, I used to be. But after having been touched by my disciples\textquotesingle{} hearts, now I only train with an altruistic spirit.\par
\medskip
And?\par
\medskip
I wish to save those youkai who are in the undeserved position at the bottom of Gensokyo\textquotesingle{}s society such as those who do not wish to fight.\par
\medskip
And what\textquotesingle{}ll happen after ya save them?\par
\medskip
I will lead them to an even higher stage.\par
\medskip
Sounds ominous to me. Just imagine a world dominated by youkai bodhisattva and youkai who have attained enlightenment. As for me, I want to create a world where there can be a pillar of support for both humans and youkai. That\textquotesingle{}s what faith is. Faith is necessary to live spiritually rich lives free of worries.\par
\medskip
That much is certain.\par
\medskip
Our methods are different, but our minds are alike.\par
\medskip
So are you holy rollers plannin\textquotesingle{} on makin\textquotesingle{} yerselves rich in the process of givin\textquotesingle{} us spiritually rich lives?\par
\medskip
Yes, but what of it? This country is full of misunderstandings about religion, with many people who oppose religion. One of the reasons for that is because religions that misuse faith are spreading also. It\textquotesingle{}s an unfortunate matter. Originally, religion was training with the goal of spiritual growth. Just like sports players use the latest theories to train their bodies to perfection, to become spiritually rich and feel happiness you need to use theories and morals, which is what religions were meant to be. Our duty is to make sure those in Gensokyo become successful yet not spiritually immature like those in the outside world.\par
\medskip
I think I get it, but hearin\textquotesingle{} that from you... Anyway, since ya brought it up, let\textquotesingle{}s change the topic to the outside world.\par
\medskip
She says \textquotedbl{}manage\textquotedbl{}, but she actually is a god residing in the shrine herself.\par
A temple with more youkai monks than human ones; the sort that could only exist in Gensokyo.\par
About a few years have passed since coming to Gensokyo.\par
In truth, everything is fine as long as the Barrier\textquotesingle{}s power remains. This is a sophism often used by youkai.\par
So that the humans could not be ruled over and used as conveniently.\par
Counting divine pillars is normally used when counting sacred entities.\par
A frog-{}associated divine entity.\par
A location in the outside world. I heard it has a sea.\par
A treasure-{}hunting companion entity.\par
Positive thinking.\par
A suspiciously manipulative way of thinking by selling things, hinging on the line \textquotedbl{}don\textquotesingle{}t you know I\textquotesingle{}m doing this for you?\textquotedbl{}.\par
\bigskip
Sumber:\par
\hspace*{1.5em}- \url{https://en.touhouwiki.net/wiki/Symposium_of_Post-mysticism/Part_1}
\end{tcolorbox}


\subsection*{Kunci}

Matriks kunci $5 \times 5$ hasil rekonstruksi:

\begin{center}
  \begin{tabular}{|c|c|c|c|c|}
    \hline
    S & Y & M & P & O \\
    \hline
    I & U & L & N & A \\
    \hline
    T & E & B & C & D \\
    \hline
    F & G & H & K & Q \\
    \hline
    R & V & W & X & Z \\
    \hline
  \end{tabular}
\end{center}

Susunan ini mengikuti pola pembentukan kunci Playfair pada umumnya: deretan kata kunci
\texttt{SYMPOIULNATE} di awal, lalu dilanjutkan sisa alfabet secara berurutan
(\texttt{BCDFGHKQRVWXZ}).

\subsection*{Referensi}

Tabel referensi frekuensi bigram Bahasa Inggris yang digunakan pada Tahap 1 mengacu pada
\cite{cornell-freq, case-ngram}. Tabel \texttt{QUADGRAM} yang dipakai sebagai fungsi objektif pada
Tahap 2 disusun dari 45 \textit{quadgram} tersering pada berkas statistik \texttt{english\_quadgrams.txt}
\cite{quadgram-data}, dengan nilai persentase dihitung terhadap total kemunculan seluruh
\textit{quadgram} pada berkas tersebut. Konsep Playfair \textit{Cipher} mengacu pada bahan kuliah
\cite{munir-klasik2}.
